\Needspace{8\baselineskip}
\section{Independent frozen-video representation control}
This additional study was specified after the original results and is post hoc. It uses the official frozen V-JEPA 2 ViT-L encoder, repository \path{facebook/vjepa2-vitl-fpc64-256}, at the following pinned commit:
\par\noindent{\small\texttt{b3c1679b7c34d3255ef3547f27c7b226aefab26f}}\par
All 24,480 original frame files across 1,980 development and 1,080 previously evaluated episodes were hash-checked. Each clip contains exactly the original eight ordered RGB frames. No additional frames, temporal resampling, or physical target is used during feature extraction. Spatial preprocessing uses the encoder's 256-pixel center crop after its native resize, and encoder evaluation uses BF16. This differs from the Qwen path's spatial processing.

The final encoder tokens are averaged to 1,024 features. Centered PCA is fitted separately on each training fold and retains 48 components; the evaluation fold never fits the projection. The final projection uses all development features only. The existing numerical head, residual target, penalties, clipping, visual shrinkage, protocol masks, generalized-bounce fallback, and covariance rule are unchanged. The same six development family folds are used. This matches residual-stage rules and output dimension, but not encoder pretraining, observation-decoder supervision, or preprocessing. The PCA map additionally stores 49,152 loadings and 1,024 centering values; these are outside the original residual-coefficient count.

\begin{table}[htbp]\centering\small
\caption{Post hoc frozen-video control. Action errors are in millimeters. Development evaluates the residual stage using the original fixed upstream models. The final column gives the existing-cohort paired difference from the original full model and an unadjusted two-sided 95\% interval.}
\begin{tabular}{lrrl}\toprule
Configuration & Development & Evaluated & Difference from full [95\% interval]\\\midrule
Numerical only & 1.371928 & 1.468770 & $+0.01565\ [-0.00053,+0.03114]$\\
Original full & 1.352231 & 1.453122 & reference\\
V-JEPA 2, PCA-48 & 1.382427 & 1.470615 & $+0.01749\ [-0.00118,+0.03560]$\\\bottomrule
\end{tabular}\end{table}

The V-JEPA-minus-numerical difference is $+0.00185$ mm with interval $[-0.00914,+0.01288]$ mm on the existing evaluation cohort. The full model has the lower point error, but neither this comparison nor V-JEPA versus full establishes a two-sided directional difference. The development V-JEPA-minus-full difference is $+0.03020$ mm, with interval $[+0.01019,+0.05058]$ mm; it does not supersede the existing-cohort uncertainty. All planned configurations are reported without changing the primary model. V-JEPA's evaluated NLL/dimension, marginal CRPS, and joint 90\% coverage are $-2.543389$, $0.01318921$, and $0.973611$, respectively. Cached embeddings, the fixed plan, PCA maps, prediction arrays, and refitting code accompany the package.

\section{Upstream lineage and an RGB-path replay}
The saved observation-decoder training records identify 328,356 parameters, 720 training episodes, and 48 old-development geometry families. The training objective combines center-of-mass pixel supervision with fixed physical-response sensitivity terms. The frozen encoder is not trained in this stage. An identifier audit finds no episode or family-identifier overlap with the later 72-family confirmation cohort. This is an identifier-level check, not an independent proof of asset or perceptual disjointness. Training receipts and source contracts are supplied; upstream training itself was not rerun.

The residual-development cohort contains these 720 upstream training episodes. Consequently, its six-fold results concern residual fitting conditional on the previously trained decoder, not out-of-fold training of every learned component. This distinction also limits interpretation of the development comparison with a video encoder that lacks this task-specific decoder supervision. The original later geometry cohort remains the primary evaluation.

The common appearance prior is a separate upstream learned component. Its recorded plan uses 12,000 first-image mappings (2,538 unique image hashes), with 8,400 training and 1,200 validation rows including repeated appearances. A frozen full Qwen pass uses the prompt ``Describe only the visible appearance of the object.'' The final-token feature undergoes float32 layer normalization, followed by a linear six-output head for three transformed means and three softplus scales. The plan specifies seeds 17, 43, and 101, 100 epochs, and minimum validation Gaussian NLL checkpoint selection; the primary path uses the frozen seed-17 checkpoint. This shared prior is also present in the numerical residual control. These are recorded training specifications, not a fresh reproduction of prior training.

An additional prior-lineage audit checks actual image-file hashes. The 9,600 fitting/selection rows contain 2,036 unique first-image hashes. Against all 1,080 later confirmation episodes (5,703 unique frame hashes), episode identifiers, family identifiers, and byte-exact image hashes have zero overlap. The saved 100-epoch history also reproduces the reported seed-17 selection at epoch 6. This does not prove perceptual or underlying asset independence, or exclude foundation-model pretraining overlap. The identifier/hash tables, training plan, and selection receipts accompany the evidence archive.

The RGB reproduction archive selects the first eight cases of the original observation index without using outcomes, covering all three protocols. Starting from 64 original RGB frames, it regenerates the first-image prior, tracking, numerical observation initialization, Qwen vision and decoder features, robust physical solves, and locked residual predictions. Prior, tracking, initialization, pooled, and position-decoder arrays reproduce exactly in the recorded environment. The maximum transformed-mean discrepancy is $5.29\times10^{-10}$, and the maximum covariance-entry discrepancy is $6.06\times10^{-14}$. Before execution, absolute tolerances were set to $10^{-6}$ for prior outputs, $10^{-4}$ for numerical inputs and decoder features, and $10^{-5}$ for posterior entries. Reference arrays are used for assertions; physical targets do not enter prediction.

The archive includes required research modules and small checkpoints, while the pinned Qwen base weights are separately downloaded and hash-checked. An independently extracted copy passes the same assertions while a Python file-read guard rejects dependencies on the original research source/data checkout. It uses the existing Python 3.11/CUDA runtime, not a newly installed environment. To preserve the original prior-head CPU batch shape, eight features are zero-padded to 256 rows; the linear head has no cross-row operation. This checks eight inference paths, not upstream training or full-dataset reproducibility. Run \path{scripts/run_access_rgb_bundle.py} in the RGB archive with its explicit model directory. Run \path{tools/score_access_video_control.py} in the Overleaf archive with Python, NumPy, and SciPy to refit the frozen-video control from bundled embeddings; its outputs are written to \path{replayed/video_control/}.
